\documentclass[11pt]{article}
\usepackage[margin=1in]{geometry}
\usepackage{booktabs}
\usepackage{amsmath}
\usepackage{graphicx}
\usepackage{hyperref}
\hypersetup{colorlinks=true, linkcolor=blue, urlcolor=blue}
\sloppy

\title{What 16{,}199 Discretionary Land-Use Decisions Look Like:\\
\large San Francisco Planning Commission, 1998--2026}
\author{Dan Post}
\date{September 2026}

\begin{document}
\maketitle

\begin{abstract}
\noindent
This memo reports what the item-level dataset contains now that it exists. Every item the San
Francisco Planning Commission heard between 1998 and mid-2026 that carries a case number ---
\textbf{16{,}199 items across 8{,}987 distinct cases} --- has been extracted into 29 structured
fields. It describes the composition of the docket, what the Commission does to what comes
before it, how long projects take, how individual commissioners vote, where in the city the
activity is, and which parts of the Planning Code are actually invoked.

\medskip\noindent
Four findings organise the rest. \textbf{Outright denial is close to non-existent} --- 1.4\%
of items across twenty-eight years, and effectively zero since 2020 --- while
\textbf{conditioning has more than tripled}, from 13.8\% of items in 1998 to 45.7\% in 2023.
The Commission has not become more permissive; it has moved from a denial technology to a
conditioning technology. \textbf{Delay is the second margin}: 45\% of all items are a repeat
appearance of a case already heard, and the 90th-percentile case now takes a year from first
hearing to last, up from five months in the late 1990s. \textbf{Discretionary review has
collapsed} as a share of the docket, from 28.6\% to 8.6\%, following the 2013 reforms.

\medskip\noindent
The memo closes on what the data cannot yet do. The regulatory burden a developer faces is
not a disposition but an \emph{expectation} formed before filing --- over delay, over
conditions, over the probability of either. The dataset now measures realised outcomes well.
Turning those into a measure of \emph{ex ante} cost and uncertainty is the modelling problem
the paper rests on, and \S\ref{sec:uncertainty} sets out four routes into it. A tested
linkage to the city's building-permit records --- \textbf{97.5\% of parsed permit numbers
match} --- shows that construction outcomes are observable on the other side of the
entitlement.
\end{abstract}

\tableofcontents
\newpage

\section{The data}

\subsection{What is in it}

One row is one item on one agenda: a request heard by the Commission on a date. The corpus
runs from January 1998 to the last scraped meeting of 2026 and holds \textbf{16{,}199 items}
carrying a case number, drawn from 985 meeting documents. Blocks with no case number ---
public comment, the Director's report, the weekly recap, adjournment --- are agenda
scaffolding and are excluded by rule.

Extraction is Claude Haiku~4.5 with six retrieved, era-matched worked examples, run over the
Batch API against a schema of 29 fields. On a frozen 139-item test half that no prompt was
tuned against, the configuration scores \textbf{96.8\%} at the field level: 97.4\% on the
1998--2014 HTML minutes and 94.0\% on the 2015--2026 PDF minutes. The companion memo,
\emph{Item-Level Field Extraction}, documents the method, the metric and the failure modes;
this memo takes the output as given and asks what it says.

\subsection{Two coverage gaps to keep in view}

\textbf{1998--2000 is thin.} Those years hold 12 meeting documents each against roughly 40 in
every year from 2001. The Commission met weekly throughout; the archive does not. Rates
computed within those years are usable, counts are not comparable, and the composition
figures should be read as describing the meetings that survive.

\textbf{2025--2026 is incomplete.} Eight documents each, against 34--41 in the preceding
years, because the scrape reached the archive before the archive reached the present. Every
count for those years is a floor, and the memo's figures mark the boundary rather than
letting a reader infer a collapse in Commission activity that did not happen.

\section{The composition of the docket}

\begin{figure}[htbp]\centering
\includegraphics[width=\textwidth]{fig_composition.pdf}
\caption{Items heard per year (top) and their composition (bottom). The thin years at each end
are archive coverage, not Commission activity.}
\end{figure}

Table~\ref{tab:composition} gives the numbers. Three movements matter.

\textbf{Discretionary review collapsed.} DR was 28.6\% of the docket in 1998 and 38.8\% in
2003 --- at its peak, the single largest category, and the mechanism by which any member of
the public could pull an as-of-right building permit in front of the Commission. By 2013 it
was 12.5\% and by 2025, 8.6\%. The timing points at the 2010--2013 DR reforms, which raised
the bar for a request to be heard at all. This is the clearest structural break in the
dataset and it changes what the Commission's docket \emph{is}: a body that spent a third of
its agenda reviewing individual permits now spends a twelfth.

\textbf{Conditional use grew to fill the space}, from 35.9\% to 48.4\% at its 2023 peak. Some
of this is composition rather than volume --- the denominator shrank as DR left --- but the
absolute count of CU items is also higher in the 2020s than the 2000s despite a smaller
docket overall.

\textbf{Legislative business roughly doubled}, from 6.6\% to 15.6\%, and \emph{informational}
items went from a rounding error to 14.8\% of the 2025 docket. The Commission increasingly
spends its time on rules rather than on projects. For a study of regulation as a market, that
is not a nuisance category: it is the supply side of rule-making appearing in the same series
as the demand side of individual permissions.

\section{What the Commission does}

\begin{figure}[htbp]\centering
\includegraphics[width=\textwidth]{fig_outcomes.pdf}
\caption{Outcome rates. Faint lines are the annual series; bold lines are a three-year centred
mean.}
\end{figure}

Table~\ref{tab:outcomes} is the central result of this memo.

\subsection{Denial is not the instrument}

Across twenty-eight years the Commission disapproved \textbf{1.4\%} of the items before it. In
1998 the rate was 4.1\%; since 2020 it rounds to zero. If regulatory stringency were measured
by refusal, San Francisco would look like one of the most permissive jurisdictions in the
country, which is not a conclusion anybody believes.

\subsection{Conditioning is}

Over the same period the share of items carrying conditions rose from \textbf{13.8\%} to
\textbf{45.7\%}. The Commission has substituted along a margin that a denial-rate measure
cannot see: it says yes, and attaches terms. Approval with conditions is the modal outcome of
a modern conditional-use hearing.

This substitution is the single most important pattern for the paper's framing. A discretionary
tax that operates through conditions rather than refusals is invisible to any measure of
approval probability, and it is precisely the kind of burden that is costly to anticipate ---
which is where uncertainty enters.

\subsection{Continuance is the third margin}

39\% of all items end in a continuance, though the rate has fallen from 46\% to 25\%.
Table~\ref{tab:outcomes} understates how much of the docket this represents, because a
continued item returns: \textbf{45\% of all 16{,}199 items are a repeat appearance of a case
already heard}. Nearly half the Commission's agenda is re-hearing.

\subsection{The disposition depends on what is asked}

Pooling all request types hides that these are different proceedings. Conditional use is
approved 54.9\% of the time and continued 35.9\%. Discretionary review is approved 1.3\% of
the time --- because ``approval'' is not its vocabulary: 26.8\% end in \emph{did not take DR}
and 15.4\% in \emph{took DR and approved}, so the modal DR outcome is the permit proceeding
essentially untouched. Variances are continued 44.6\% of the time and approved 15.4\%, which
reflects that the Zoning Administrator, not the Commission, is the usual variance authority.

\section{Delay}

\begin{figure}[htbp]\centering
\includegraphics[width=\textwidth]{fig_delay.pdf}
\caption{Repeat hearings and elapsed time, by year of first hearing.}
\end{figure}

A case heard more than once is a project the Commission saw and did not finish with. Chaining
items by case number gives 8{,}987 cases, of which \textbf{37.9\% were heard more than once};
the record is a case that took 3{,}423 days --- nine years and four months --- across its
appearances.

Table~\ref{tab:delay} shows the distribution by cohort. The median has been stable at about
five weeks, which is roughly the Commission's own calendar cycle: a routine continuance moves
an item two or three meetings. \textbf{The tail is what moved.} The 90th percentile went from
147 days for cases first heard in the late 1990s to \textbf{360 days} for the 2015--2019
cohort. The typical case is not slower. The unlucky case is more than twice as slow.

That is a statement about \emph{variance}, and it matters more than the median for a developer
deciding whether to file. Elapsed time also sorts on outcome in a way that is intuitive but
worth having measured: cases ending in approval run a median 28 days from first to last
hearing, those ending \emph{continued indefinitely} run 84, and those ending in no action run
116.

\medskip\noindent\textbf{Two cautions.} First, this measures hearing-to-hearing time, not the
full entitlement clock: it begins when an item first reaches an agenda, which is already well
downstream of filing. Second, recent cohorts are right-censored --- a case first heard in 2024
cannot yet show a three-year tail --- so the recent decline in Table~\ref{tab:delay} is partly
an artefact and is flagged as such rather than read as an improvement.

\section{Commissioners}

\begin{figure}[htbp]\centering
\includegraphics[width=0.95\textwidth]{fig_commissioners.pdf}
\caption{Dissent rate for commissioners with 200 or more recorded votes.}
\end{figure}

The minutes record roll-call votes on both sides, so this is a census rather than a sample:
\textbf{79{,}855 individual votes by 84 named commissioners}. Table~\ref{tab:commissioners}
gives the twenty most active.

The striking feature is how \textbf{rarely anyone dissents}. Rates run from 1.1\% to 7.9\%, and
the Commission is close to unanimous on the overwhelming majority of items. The spread is
nonetheless real and persistent: Moore (7{,}608 votes, 7.2\%), Sugaya (7.9\%) and Olague
(7.8\%) dissent five to seven times as often as Diamond (1.1\%), Borden (1.2\%) or Koppel
(1.4\%), and these are not small samples.

Two things this enables. First, a commissioner fixed effect is estimable for the two dozen
members with enough votes, which makes it possible to ask whether \emph{who is sitting} moves
outcomes conditional on what is asked. Second, because the roll call is per item and the
composition of the Commission turns over, the panel has within-project variation in
composition for cases heard repeatedly --- the same project, a different room.

\section{Geography}

\begin{figure}[htbp]\centering
\includegraphics[width=\textwidth]{fig_geography.pdf}
\caption{Assessor blocks appearing most often before the Commission, 1998--2026.}
\end{figure}

\textbf{84.1\% of items carry a usable parcel key} --- an assessor block and at least one lot
--- covering 2{,}557 distinct blocks and 8{,}295 distinct parcels. This is the join that makes
everything downstream possible, and it is the reason \texttt{project\_address} was demoted to
a validation field: the block-and-lot pair is what the city's own records key on, and it
scores 98--100\% where the address needs a geocoder.

Concentration is high. The most-heard block (3736, South of Market at the Rincon Hill/Transbay
edge) appears 95 times; the top ten are almost entirely SoMa and the Downtown core --- 3704,
3721, 3753, 3787, 3788 --- with two North Beach blocks (131, 287) as the exception. These are
the blocks where large projects, phased approvals and repeated returns concentrate.

Hot spots are a genuine finding but they need care: a block appears often both because a lot
of building is proposed there \emph{and} because a single large project can generate a dozen
items across a decade. Separating those is a matter of collapsing to the case level and is
straightforward with the chain construction of \S4.

\section{Which parts of the Planning Code are actually used}

Table~\ref{tab:citations} lists the sections cited in the request text. \textbf{1{,}347
distinct sections} appear at least once; the distribution is extremely skewed.

The result is a validity check as much as a finding. Nothing told the extractor what the
Planning Code contains, yet the mapping from section to request type recovers the Code's own
structure: \S303 (conditional use procedure) is cited by 2{,}468 conditional-use items;
\S309 (permit review in C-3 districts) by 252 downtown-project items; \S329 (large project
authorization) by 179 LPA items; \S321 (office development annual limit) by 209
office-allocation items. Where the extraction assigns a request type, the cited section
independently agrees.

Substantively, four sections carry most of the discretionary weight:

\begin{itemize}
\item \textbf{\S31.04} (2{,}899 items) is the Administrative Code's approval-action provision
for CEQA. Its ubiquity is procedural boilerplate, but it dates the moment at which the
entitlement becomes appealable.
\item \textbf{\S303} (2{,}701) is the conditional-use standard --- the ``necessary or
desirable'' test, which is the discretion.
\item \textbf{\S101.1} (923) is the eight Priority Policies. It appears on legislative items
and is the general-welfare hook that makes almost any objection cognisable.
\item \textbf{\S134} (912, rear yard) and \textbf{\S140} (358, dwelling-unit exposure) are the
two bulk standards that generate the most variance requests --- the physical constraints that
bind hardest on infill.
\end{itemize}

\S317, the demolition control, appears on 660 items and is the single most direct instrument
of housing-stock protection in the dataset.

\section{What the patterns say}

\subsection{Stringency moved from the extensive to the intensive margin}

The headline pattern is a substitution: denial fell to zero while conditioning tripled. A
regulator that rarely refuses but usually attaches terms is not a permissive regulator; it is
one whose burden is expressed in the content of the approval rather than in its existence.
Any measure of regulatory stringency built on approval rates will read San Francisco as
permissive and will be wrong.

\subsection{The binding constraint is time, and specifically its variance}

Half the docket is re-hearing. The median case moves at the speed of the meeting calendar, but
the 90th percentile more than doubled. For a developer with financing costs, the relevant
quantity is not expected delay but the width of its distribution --- and that width has grown
even as the centre has not.

\subsection{Discretionary review's collapse is a natural experiment}

DR fell from 38.8\% of the docket to 8.6\% following reforms that changed who could trigger a
hearing. That is a sharp, dated, policy-induced change in the \emph{intensity} of discretionary
review --- exactly the parameter the supply model needs to move --- and it is identified by
something other than the demand for housing.

\section{From this dataset to a cost of uncertainty}
\label{sec:uncertainty}

The honest position is that this dataset measures \textbf{realised outcomes} well and
\textbf{ex ante burden} not at all. A developer choosing whether to file faces a distribution
over delay, over conditions and over approval; the dataset records one draw from it. Four
routes seem worth trying, in increasing order of ambition.

\subsection{The problem with \texttt{conditions\_imposed} as it stands}

A concrete obstacle first, because it is the binding one. \texttt{conditions\_imposed} is
populated on 25.9\% of items --- but its median length is \textbf{three characters}. It is a
yes/no flag. The minutes very often say that conditions were imposed without enumerating them,
and the schema faithfully records that. \texttt{modifications} does carry text (median 108
characters) but is populated on only 10.7\% of items.

So the dataset can currently answer \emph{whether} a project was conditioned, and almost never
\emph{how}. That is the gap between the pattern in \S3 and a cost measure, and closing it is
the highest-value next step. Three options:

\begin{enumerate}
\item \textbf{Motions and resolutions.} The Commission's conditions live in the numbered motion
or resolution, not in the minutes. \texttt{action\_instrument} and
\texttt{action\_instrument\_no} are extracted on 40\% of items and give the document number
directly. Those documents are public and are the actual text of the burden.
\item \textbf{The Planning Department's own condition-of-approval templates.} Conditions are
substantially standardised. A modest taxonomy --- affordability, transportation demand
management, open space, hours of operation, design review, monitoring --- would convert most
condition text into a handful of indicators.
\item \textbf{Cost-relevant coding.} Not every condition is costly. The distinction that
matters is between conditions that change the project (units, height, envelope) and those that
change its obligations (fees, monitoring, reporting). The \texttt{project\_modified} flag
already separates these in principle.
\end{enumerate}

\subsection{Route one: realised delay as a predictable quantity}

The chain construction gives, for every case, the number of hearings and elapsed days. Regress
that on what was observable \emph{at first hearing} --- request type, district, block, staff
recommendation, speaker turnout --- and the fitted value is an expected delay conditional on
observables, while the residual variance is the part a developer cannot forecast. That
residual is the natural first proxy for uncertainty, and it is available now without new data.

\subsection{Route two: the staff recommendation as a prior}

\texttt{preliminary\_recommendation\_category} is populated on 78.7\% of items and is issued
\emph{before} the hearing. It is a public, dated signal of the expected outcome. Two
quantities follow: the follow rate (how often the Commission does what staff asked), which
prices how much residual discretion the Commission exercises; and the subset where the
Commission departs from staff, which is where the unpredictable burden concentrates. A model
in which the developer's prior is the staff recommendation and the Commission's deviation is
the shock is both estimable and economically legible.

\subsection{Route three: speakers as a mobilisation shock}

Public turnout is recorded per item with a stance where the minutes mark one. Turnout is
plausibly the most visible predictor of trouble, and it is observable to the developer only
imperfectly at filing. Its usefulness is limited by the field's own accuracy --- 69.6\% on
the test half, the weakest field in the schema --- so this route needs the stance work
finished before it can carry weight.

\subsection{Route four: the option value of not filing}

The deepest version of the question is not measurable in this dataset at all, and it should be
said plainly. Projects that were never filed because the process was too uncertain do not
appear in Commission minutes. Any estimate of the cost of uncertainty built only on the
observed docket is conditioned on selection into it. The permit records of \S9 are the only
route to the counterfactual: parcels where a permit was filed and abandoned, or where an
entitlement was granted and never built, are the observable trace of a project that lost its
bet.

\section{Linking to the city's permit records}
\label{sec:permits}

The dataset's value multiplies if the entitlement can be followed to construction. This
section reports a first, tested pass.

\subsection{Permit numbers parse and they match}

3{,}932 items mention a building or demolition permit in their request text. A permit number
can be parsed from \textbf{2{,}778} of them, yielding 1{,}752 distinct permits in three printed
forms --- the modern \texttt{2013.10.21.9832}, the packed \texttt{201310219832}, and the
1990s eight-digit \texttt{9725973}.

Tested against DataSF's Building Permits dataset (1{,}295{,}048 rows, dataset
\texttt{i98e-djp9}) on a random sample of 120 parsed numbers, \textbf{117 matched --- 97.5\%}.
The join works, and it works on the exact identifier the minutes print.

\subsection{What the match buys}

The DBI record carries \texttt{filed\_date}, \texttt{issued\_date},
\texttt{first\_construction\_document\_date}, \texttt{status}, \texttt{estimated\_cost},
\texttt{revised\_cost}, \texttt{proposed\_units} and \texttt{existing\_units}. So for a matched
item we can observe, on the far side of the discretionary hearing:

\begin{itemize}
\item \textbf{whether the permit was ever issued}, and how long after entitlement;
\item \textbf{whether construction started}, via the first construction document;
\item \textbf{whether it was withdrawn or expired} --- the abandonment outcome that \S8.5
identifies as the trace of a lost bet;
\item \textbf{how much it cost} and \textbf{how many units} resulted, against how many were
proposed at the hearing.
\end{itemize}

Sampled records show all four states occurring: permits complete, permits issued but not
completed, and permits withdrawn after filing.

\subsection{The parcel route is broader but noisier}

Where no permit number is printed, the block-and-lot key reaches further: on a random sample
of 150 parcels from the minutes, \textbf{68\% had at least one DBI permit}. But the median
matched parcel carries \emph{15} permits, most of them plumbing, electrical and small
alterations irrelevant to the entitlement. The parcel join therefore needs a date window and a
permit-type filter before it identifies the right permit, and its precision should be
validated against the permit-number subset --- which is exactly what that subset is for.

\subsection{What is not yet established}

Whether the \emph{conditions} attached at the hearing are observable in the permit record is
unresolved. DBI records the permit's own status and scope, not the Planning Commission's
conditions of approval. Recovering those means the motions and resolutions of \S8.1, not
DataSF. The permit join gives outcomes; it does not give the burden.

\section{Where the data lives}

The corpus pass is stored under \verb|$MFHR_DATA_ROOT/extraction/corpus_v2_g3/| in three forms
kept separately and never conflated:

\begin{itemize}
\item \texttt{raw/} --- the model's replies exactly as returned, never rewritten;
\item \texttt{interim/} --- unwrapped from the evidence envelope and type-coerced;
\item \texttt{clean/} --- normalised; the analysis input for this memo;
\item \texttt{evidence/} --- the supporting spans, and \texttt{verification\_failures.csv};
\item \texttt{manifest.json} --- batch ids, run ids, token counts, cost, prompt SHA, schema and
gold versions.
\end{itemize}

The separation is not fastidiousness. Every serious defect found while building the gold set
was a \emph{normalisation} defect --- one silently emptied a field on 122 of 232 gold records
--- and each was recoverable only because the untouched replies had been kept. Re-deriving
\texttt{clean} from \texttt{raw} is a local loop; re-querying the corpus is \$69.

\section{Limitations}

\begin{enumerate}
\item \textbf{Coverage is uneven at both ends.} 1998--2000 and 2025--2026 hold roughly a fifth
of the meeting documents of the intervening years. Rates are usable; counts are not.
\item \textbf{Field accuracy is 96.8\%, not 100\%.} At the item level the errors compound: an
analysis conditioning on five fields at once is working with materially fewer clean rows than
the field-level rate suggests, and whole-record accuracy should be reported wherever that is
the operative unit.
\item \textbf{\texttt{conditions\_imposed} is a flag, not a description}, which is the binding
limitation on measuring stringency and the subject of \S8.1.
\item \textbf{Case-number chains are an imperfect project identifier.} A project can change
case number across a re-filing, and one case number can cover several agenda sub-items heard
the same day. The chains in \S4 are a lower bound on how often projects return.
\item \textbf{Delay is measured from first hearing}, which is already downstream of filing.
The pre-hearing clock --- often the longer one --- is invisible here and lives in the
Department's own records.
\item \textbf{Speaker stance is the weakest field} (69.6\%) and the stance rule changed after
the scored run. Any mobilisation analysis should wait for a re-run.
\item \textbf{No causal claim is made anywhere in this memo.} The DR collapse in \S7.3 is
described as a candidate natural experiment, not as an estimate.
\end{enumerate}

% GENERATED BY analyze_corpus.py — do not edit by hand.
\begin{table}[htbp]\centering
\caption{Composition of the Commission's item-level docket, per cent of items heard that year. Families group the 17 \texttt{request\_type} values; `Legislative' is Planning Code, map and General Plan amendments.}\label{tab:composition}
\resizebox{\textwidth}{!}{%
\begin{tabular}{lrrrrrrrrrr}\toprule
Year & Conditional use & Discretionary review & Variance & Legislative & Large project /\allowbreak  downtown & CEQA /\allowbreak  appeals & Historic & Informational & Other & Items\\\midrule
1998 & 35.9 & 28.6 & 3.9 & 6.6 & 5.6 & 14.9 & 2.9 & 0.7 & 1.0 & 715\\
2003 & 33.8 & 38.8 & 9.1 & 6.8 & 2.1 & 6.4 & 1.0 & 0.8 & 1.4 & 809\\
2008 & 27.1 & 24.9 & 8.5 & 17.7 & 5.6 & 8.0 & 1.4 & 4.3 & 2.5 & 851\\
2013 & 42.3 & 12.5 & 8.0 & 15.7 & 8.0 & 7.4 & 0.2 & 4.1 & 1.8 & 513\\
2018 & 37.6 & 19.6 & 7.6 & 15.6 & 5.0 & 4.5 & 0.4 & 5.2 & 4.5 & 540\\
2023 & 48.4 & 14.4 & 2.9 & 15.6 & 5.8 & 3.1 & 0.9 & 4.4 & 4.4 & 455\\
\bottomrule\end{tabular}}\end{table}

\begin{table}[htbp]\centering
\caption{Commission outcomes, per cent of items heard that year. `Conditions' and `modified' are the orthogonal flags, not dispositions: a project can be approved with conditions and modified at once.}\label{tab:outcomes}
\begin{tabular}{lrrrrr}\toprule
Year & Continued & Conditions & Modified & Disapproved & Took DR\\\midrule
1998 & 46.4 & 13.8 & 9.2 & 4.1 & 5.4\\
2003 & 48.2 & 27.1 & 11.2 & 0.6 & 5.5\\
2008 & 44.5 & 20.7 & 8.6 & 1.7 & 5.0\\
2013 & 27.2 & 40.0 & 7.4 & 1.6 & 3.0\\
2018 & 35.9 & 29.7 & 9.0 & 1.5 & 5.0\\
2023 & 24.6 & 45.7 & 7.5 & 0.0 & 3.8\\
\midrule All years & \textbf{39.0} & \textbf{26.4} & \textbf{8.8} & \textbf{1.4} & \textbf{5.0}\\
\bottomrule\end{tabular}\end{table}

\begin{table}[htbp]\centering
\caption{Elapsed time from first to last hearing, for cases heard more than once, by five-year cohort of first hearing. A case is a \texttt{case\_number}; the corpus ends mid-2025, so recent cohorts are right-censored and their tails are understated.}\label{tab:delay}
\begin{tabular}{lrrrr}\toprule
First heard & Cases & Median days & 90th pct & Max\\\midrule
1995--1999 & 314 & 35 & 147 & 665\\
2000--2004 & 845 & 35 & 151 & 3157\\
2005--2009 & 798 & 35 & 163 & 2408\\
2010--2014 & 553 & 35 & 209 & 2737\\
2015--2019 & 457 & 63 & 360 & 3423\\
2020--2024 & 429 & 49 & 309 & 1316\\
\midrule All & 3,409 & 35 & 210 & 3423\\
\bottomrule\end{tabular}\end{table}

\begin{table}[htbp]\centering
\caption{Commissioners with 500 or more recorded votes. Dissent is a vote recorded in the NOES column; the Commission votes by roll call and the minutes record both sides, so this is a census of recorded votes rather than a sample.}\label{tab:commissioners}
\begin{tabular}{lrrrl}\toprule
Commissioner & Votes & Noes & Dissent rate & Years\\\midrule
Moore & 7,608 & 550 & 7.2\% & 2006--2026\\
Antonini & 7,012 & 264 & 3.8\% & 2002--2016\\
Olague & 3,975 & 309 & 7.8\% & 2004--2012\\
Sugaya & 3,734 & 295 & 7.9\% & 2006--2014\\
S. Lee & 3,147 & 80 & 2.5\% & 2002--2008\\
Koppel & 3,099 & 43 & 1.4\% & 2016--2024\\
W. Lee & 2,910 & 49 & 1.7\% & 2002--2008\\
Fong & 2,665 & 48 & 1.8\% & 2010--2019\\
Joe & 2,560 & 145 & 5.7\% & 1998--2002\\
Chinchilla & 2,451 & 81 & 3.3\% & 1998--2002\\
Borden & 2,426 & 30 & 1.2\% & 2008--2014\\
Theoharis & 2,403 & 48 & 2.0\% & 1998--2002\\
Hillis & 2,224 & 44 & 2.0\% & 2012--2019\\
Hughes & 2,175 & 56 & 2.6\% & 2002--2006\\
Miguel & 1,964 & 32 & 1.6\% & 2008--2012\\
Imperial & 1,963 & 141 & 7.2\% & 2020--2026\\
Diamond & 1,850 & 21 & 1.1\% & 2019--2024\\
Johnson & 1,829 & 28 & 1.5\% & 2014--2020\\
Richards & 1,613 & 51 & 3.2\% & 2014--2019\\
Mills & 1,611 & 20 & 1.2\% & 1998--2000\\
\bottomrule\end{tabular}\end{table}

\begin{table}[htbp]\centering
\caption{Planning Code sections cited in the request text, with the two request types that cite them most. The mapping recovers the Code's own structure without being told it, which is a validity check on the extraction as much as a finding.}\label{tab:citations}
\resizebox{\textwidth}{!}{%
\begin{tabular}{lrl}\toprule
Section & Items & Cited mainly by\\\midrule
\S31.04 & 2,899 & \texttt{conditional\_use} (1,652), \texttt{discretionary\_review} (858)\\
\S303 & 2,701 & \texttt{conditional\_use} (2,468), \texttt{conditional\_use\_modification} (141)\\
\S101.1 & 923 & \texttt{planning\_code\_amendment} (707), \texttt{general\_plan\_amendment} (90)\\
\S134 & 912 & \texttt{variance} (519), \texttt{conditional\_use} (190)\\
\S302 & 726 & \texttt{planning\_code\_amendment} (586), \texttt{rezoning\_map\_amendment} (103)\\
\S317 & 660 & \texttt{conditional\_use} (397), \texttt{discretionary\_review} (215)\\
\S309 & 454 & \texttt{downtown\_project} (252), \texttt{large\_project\_authorization} (89)\\
\S140 & 358 & \texttt{variance} (166), \texttt{conditional\_use} (81)\\
\S209.1 & 307 & \texttt{conditional\_use} (288), \texttt{planning\_code\_amendment} (10)\\
\S321 & 252 & \texttt{office\_allocation} (209), \texttt{large\_project\_authorization} (13)\\
\S209.3 & 247 & \texttt{conditional\_use} (227), \texttt{conditional\_use\_modification} (10)\\
\S304 & 231 & \texttt{conditional\_use} (191), \texttt{conditional\_use\_modification} (21)\\
\S135 & 214 & \texttt{variance} (86), \texttt{conditional\_use} (46)\\
\S151 & 202 & \texttt{variance} (106), \texttt{conditional\_use} (61)\\
\S329 & 189 & \texttt{large\_project\_authorization} (179), \texttt{planning\_code\_amendment} (7)\\
\S121.2 & 188 & \texttt{conditional\_use} (167), \texttt{planning\_code\_amendment} (16)\\
\S161 & 181 & \texttt{conditional\_use} (140), \texttt{downtown\_project} (13)\\
\S121.1 & 180 & \texttt{conditional\_use} (168), \texttt{planning\_code\_amendment} (7)\\
\bottomrule\end{tabular}}\end{table}

\begin{table}[htbp]\centering
\caption{What the item table can be joined to, and how far it reaches.}\label{tab:linkage}
\begin{tabular}{lrl}\toprule
Join key & Coverage & Reaches\\\midrule
Assessor block & 84.5\% & 2,557 distinct blocks\\
Block $+$ lot (APN) & 84.1\% & 8,295 distinct parcels\\
Building-permit number & 17.1\% & 1,752 distinct permits, 98\% matched in DBI\\
\bottomrule\end{tabular}\end{table}


\end{document}
